% ── CAPÍTULO 7: EJEMPLOS Y TUTORIALES ───────────────────────
\section{Ejemplos y Tutoriales}
\label{sec:ejemplos}

\subsection{Ejemplo 1: Red N5B con k=3 (Básico)}

Este tutorial muestra cómo analizar la red de 5 nodos N5B con una 3-partición.

\textbf{Paso 1 — Verificar el archivo de entrada:}

\begin{lstlisting}[style=shell]
# Verificar que el CSV existe
ls src/shared/input/N5B.csv
# Debe mostrar el archivo (32 filas de datos + encabezado)
\end{lstlisting}

\textbf{Paso 2 — Ejecutar KGeoMIP con k=3:}

\begin{lstlisting}[style=python, caption={Ejecutar KGeoMIP sobre N5B con 3 grupos.}]
# Desde la raiz del proyecto:
from src.strategies.GeoMIP.kgeomip import KGeoMIP
from src.models.core.system import System

sistema = System.desde_csv("src/shared/input/N5B.csv")
kgeo = KGeoMIP(sistema)
resultado = kgeo.ejecutar_k(k=3)

print(f"Perdida (EMD): {resultado.perdida:.4f}")
print(f"Particion: {resultado.particion}")
print(f"Tiempo: {resultado.tiempo:.3f} s")
\end{lstlisting}

\textbf{Salida esperada:}

\begin{lstlisting}[caption={Salida del ejemplo 1.}]
Perdida (EMD): 0.1180
Particion: [[(1,0)], [(1,1),(0,0)], [(1,2),(0,1),(0,2)]]
Tiempo: 0.021 s
\end{lstlisting}

\textbf{Interpretación:} La 3-partición divide los 5 nodos en tres grupos. La pérdida de 0.118 es aproximadamente la mitad que la bi-partición (0.214), lo que indica que los tres grupos explican mejor la dinámica del sistema de forma independiente.

\subsection{Ejemplo 2: Red N10A con k=4 (Intermedio)}

\textbf{Paso 1 — Ejecutar el benchmark selectivo:}

\begin{lstlisting}[style=python, caption={Benchmark selectivo para N10A.}]
import time
from src.strategies.QNodes.kqnodes import KQNodes
from src.models.core.system import System

sistema = System.desde_csv("src/shared/input/N10A.csv")
kqn = KQNodes(sistema, k=4)

inicio = time.time()
resultado = kqn.resolver()
fin = time.time()

print(f"Modo: {'exacto' if resultado.modo == 'exacto' else 'greedy'}")
print(f"Perdida: {resultado.perdida:.4f}")
print(f"Tiempo total: {fin - inicio:.2f} s")
\end{lstlisting}

\textbf{Resultado esperado:} El sistema reportará modo \textit{greedy} (porque $S(10,4) = 34{,}105 > 10{,}000$) con una pérdida aproximada de 0.148.

\subsection{Ejemplo 3: Comparar GeoMIP vs KGeoMIP}

\begin{lstlisting}[style=python, caption={Comparación directa entre GeoMIP y KGeoMIP(k=2).}]
from src.strategies.GeoMIP.geomip  import GeoMIP
from src.strategies.GeoMIP.kgeomip import KGeoMIP
from src.models.core.system import System

sistema = System.desde_csv("src/shared/input/N5B.csv")

geo  = GeoMIP(sistema)
kgeo = KGeoMIP(sistema)

r_geo  = geo.resolver()
r_kgeo = kgeo.ejecutar_k(k=2)

dif = abs(r_geo.perdida - r_kgeo.perdida)
print(f"GeoMIP  loss: {r_geo.perdida:.6f}")
print(f"KGeoMIP loss: {r_kgeo.perdida:.6f}")
print(f"Diferencia:   {dif:.2e}  ({'OK' if dif < 1e-5 else 'FALLO'})")
\end{lstlisting}

La diferencia debe ser menor que $10^{-5}$, validando TV-02.

\subsection{Ejemplo 4: Ejecutar el Benchmark Completo y Generar Dashboard}

\begin{lstlisting}[style=shell, caption={Flujo completo: benchmark + dashboard.}]
# Ejecutar todas las estrategias
python src/shared/run/benchmark.py

# Generar el dashboard interactivo
python src/shared/run/dashboard.py

# Abrir el dashboard (Windows)
start dashboard.html

# Abrir el dashboard (Linux)
xdg-open dashboard.html

# Abrir el dashboard (macOS)
open dashboard.html
\end{lstlisting}

\begin{tip}
El dashboard permite hacer zoom, exportar gráficas como PNG y comparar estrategias de forma interactiva. Use la barra de herramientas de Plotly en la esquina superior derecha de cada gráfica.
\end{tip}
